% SPDX-License-Identifier: CC-BY-4.0
% Copyright (c) 2026 Martin Schröder <info@swedishembedded.com>

\section{Limitations}\label{sec:limitations}

The structure claims show that definitions, constants, tests and messages exist at the pinned
commits; they do not show that the code behaves accordingly. Behaviour is established only by
measurement, and the only measurements here are those of the training substrate
(\cref{sec:substrate}): they show that steps fit and run, not that the adapter changes
generation. The qualification of the 27B model used random weights for the long sequences and a
handful of examples for the real checkpoint, and no sequence longer than 16384 tokens was run.

The requirements of \cref{sec:requirements} that are marked planned are not implemented: the
retention check does not test acquired facts, a workspace pass is an exit status, and the
practice gate has no canaries. The mechanisms of corpus expansion beyond the first pass
(variants, retention queue, rehearsal, claims) are described from the code and none was
evaluated for the value it adds.

\paragraph{Single greedy runs and run-to-run variation.} The development probes of \cref{sec:pending} (the supervision ablation, the recovery dose-response and the generation-2 rehearsal ablation) are one run per arm at greedy decoding. The same adapter evaluated three times on the same benchmark decoded 30, 31 and 26 of 46 tasks, and the three reports' verdicts differ pairwise on 9, 12 and 11 of the 46 tasks; a single greedy arm therefore carries about 2.5 tasks of standard deviation, so the paired p-values of single-run comparisons overstate their precision.\cl{REP1} The variation is localised, and it is not run-to-run randomness: it is the co-batch.\cl{REP2} A single-row decode is reproducible. Five sequential identical task-derived requests (the real coding system prompt, the first user request and eight tool schemas from a stored record), sent greedily with the requested seed echoed by the server, returned identical completion token ids and bitwise identical per-token log-probabilities; six identical requests sent at once returned two distinct token sequences, first differing at token seven, and eight at once also two, while two at once were identical. The arithmetic a batch selects explains it: the int8 projections run a block that is bit-identical below five rows and switch to an int8 tensor-core block from five rows, and the fused decode attention sizes its split-key count from the batch, so a row's result depends on which sequences share its batch. An arm measured with eight requests in flight therefore draws a fresh batch composition on every run, which is what the three greedy evaluations above differ by. Evaluation at one request in flight is deterministic and is the instrument any single comparison should be checked against; the registered multi-seed protocol averages the rest, and an exact batched decode - grouped int8 projections and a split count pinned to the context rather than to the batch - is the change that would remove the variation at full throughput.

\paragraph{The coding-task experiment.} Its claim is not measured. What has run is the
teacher's first pass, an underpowered pilot of the headline that does not show it, two CE8
training runs that the registered memorisation guard stopped before evaluation, and a pilot of
CE8-T and its shuffled-instruction control on the first qualification half that does not show an
improvement either (\cref{sec:pending}); everything else said about it in \cref{sec:protocol,sec:pending,sec:replicate}
is a registered rule or a statement about code, and the sign-test table is a design calculation.
The guard as registered reads one step's training loss and any rise at the last evaluation, so it
can stop a run on noise; it was kept unchanged.\cl{M48,M49} Both pilots of the adapter are inside the same
variation the base shows against itself on a second protocol seed.\cl{M57} The registration fixes one model size for the first
training result (Qwen3-8B, because training, saving and serving an adapter work end to end for
it), trained on a stronger teacher's verified trajectories because the untrained 8B produced
too few edits to supply its own (so a gain is distillation from the teacher, not
self-improvement; the self-improving 27B is the separate experiment CE8),\cl{PR15,PR16} one agent harness, and runs with process isolation \code{none} on the node (a solver's commands are
confined by the kernel to their tree and have no network, but the node is not a container;
\cref{sec:audit}).\cl{AU09,CN01,CN08} Mutation-seeded
tasks test whether a model can undo a synthetic fault, not whether it can debug a historical one;
tasks from repository history are registered as a source, not yet built.\cl{PR10} A family is the file a
fault was seeded in for the authored repositories, so clusters are files, not subsystems.\cl{PR02}
An exam of one repository may hold too few non-tied families to decide, and the registration
labels such a comparison a pilot instead of changing the rule. The base model may have seen
public repositories in pretraining; the paired base arm bounds that effect, it does not remove it.
Agent rollouts of realistic length make training cost unmeasured for the 8B model; the only
training rate measured here is that of the 27B architecture (\cref{sec:throughput}), on one card,
one revision and one sequence length.
